% GENERATED FILE. Edit canonical lessons, then run npm run generate:books.
% canonical-source-hash: fnv1a64:23e54d80edb64324
% canonical-lessons: GU-C333-pari, GU-C333-jadu, GU-C333-camatkar, GU-C333-gayab, GU-C333-dekhavu

\chapter{Fairies and Magic}
\label{ch:gu-fairies-and-magic}
\hlchaptermodality{333}

\begin{chapteropening}
\textbf{By the end of this chapter:} \emph{I can tell what magic happened in a story.}

Complete the last lesson of chapter 333: I can tell what magic happened in a story.
\end{chapteropening}

\section[\texorpdfstring{\gu{પરી} (parī)}{પરી (parī)}]{\gu{પરી} (parī) — a fairy}

\label{lesson:GU-C333-pari}



\begin{quote}
Before the new one: say the Gujarati for a goddess, then the Gujarati for a boon.
\end{quote}

\subsection*{You'll want to know: \gu{પરી}}
\textbf{\gu{પરી}} — \emph{parī} — \textquotedblleft{}a fairy\textquotedblright{}.

\begin{quote}
\gu{વાર્તામાં} \gu{એક} \gu{સુંદર} \gu{પરી} \gu{હતી}.
\end{quote}

\begin{quote}
\emph{vārtāmā̃ ek sundar parī hatī.}
\end{quote}

\begin{quote}
There was a beautiful fairy in the story.
\end{quote}

A feminine noun.

\begin{cousinweb}
\textbf{\gu{પરી}} (\emph{parī}) is usually traced to Persian \emph{parī}, \textquotedblleft{}a fairy\textquotedblright{}, the winged spirit of Persian tales. It is feminine: \textbf{\gu{સુંદર} \gu{પરી}}.
\end{cousinweb}

\begin{grammarlens}[title={What you've built}]
One more word for telling what happened.
\end{grammarlens}

\subsection*{Your turn}
\begin{itemize}
\raggedright
  \item \emph{Say it:} \emph{parī}
  \item \emph{Say it:} \emph{parī}, once more
  \item \emph{Say it:} \emph{vārtāmā̃ ek sundar parī hatī.} — There was a beautiful fairy in the story.
  \item \emph{Recall:} say the Gujarati for a goddess, then the Gujarati for a boon, then say \emph{parī} again
\end{itemize}

\begin{tcolorbox}[breakable,colback=teal!4,colframe=teal!35!black,title={Before you move on}]
What does \gu{પરી} mean? (\textquotedblleft{}A fairy\textquotedblright{}.) Say it once more.
\end{tcolorbox}
\section[\texorpdfstring{\gu{જાદુ} (jādu)}{જાદુ (jādu)}]{\gu{જાદુ} (jādu) — magic}

\label{lesson:GU-C333-jadu}



\begin{quote}
Before the new one: say the Gujarati for a boon, then the Gujarati for a fairy.
\end{quote}

\subsection*{You'll want to know: \gu{જાદુ}}
\textbf{\gu{જાદુ}} — \emph{jādu} — \textquotedblleft{}magic\textquotedblright{}.

\begin{quote}
\gu{પરીના} \gu{જાદુથી} \gu{ઝૂંપડું} \gu{મહેલ} \gu{બની} \gu{ગયું}.
\end{quote}

\begin{quote}
\emph{parīnā jāduthī jhū̃pḍũ mahel banī gayũ.}
\end{quote}

\begin{quote}
By the fairy's magic the hut turned into a palace.
\end{quote}

A masculine noun.

\begin{cousinweb}
\textbf{\gu{જાદુ}} (\emph{jādu}) is usually traced to Persian \emph{jādū}, \textquotedblleft{}magic\textquotedblright{}. It is masculine: \textbf{\gu{પરીનો} \gu{જાદુ}}, the fairy's magic; \textbf{\gu{જાદુથી}}, by magic.
\end{cousinweb}

\begin{grammarlens}[title={What you've built}]
One more word for telling what happened.
\end{grammarlens}

\subsection*{Your turn}
\begin{itemize}
\raggedright
  \item \emph{Say it:} \emph{jādu}
  \item \emph{Say it:} \emph{jādu}, once more
  \item \emph{Say it:} \emph{parīnā jāduthī jhū̃pḍũ mahel banī gayũ.} — By the fairy's magic the hut turned into a palace.
  \item \emph{Recall:} say the Gujarati for a boon, then the Gujarati for a fairy, then say \emph{jādu} again
\end{itemize}

\begin{tcolorbox}[breakable,colback=teal!4,colframe=teal!35!black,title={Before you move on}]
What does \gu{જાદુ} mean? (\textquotedblleft{}Magic\textquotedblright{}.) Say it once more.
\end{tcolorbox}
\section[\texorpdfstring{\gu{ચમત્કાર} (camatkār)}{ચમત્કાર (camatkār)}]{\gu{ચમત્કાર} (camatkār) — a miracle}

\label{lesson:GU-C333-camatkar}



\begin{quote}
Before the new one: say the Gujarati for a fairy, then the Gujarati for magic.
\end{quote}

\subsection*{You'll want to know: \gu{ચમત્કાર}}
\textbf{\gu{ચમત્કાર}} — \emph{camatkār} — \textquotedblleft{}a miracle\textquotedblright{}.

\begin{quote}
\gu{બધાએ} \gu{કહ્યું} \gu{કે} \gu{આ} \gu{ચમત્કાર} \gu{છે}.
\end{quote}

\begin{quote}
\emph{badhāe kahyũ ke ā camatkār che.}
\end{quote}

\begin{quote}
Everyone said it was a miracle.
\end{quote}

A masculine noun.

\begin{cousinweb}
\textbf{\gu{ચમત્કાર}} (\emph{camatkār}) is a tatsama, Sanskrit \emph{camatkāra}, \textquotedblleft{}astonishment, something wonderful\textquotedblright{} (Monier-Williams): what makes people cry out in wonder.
\end{cousinweb}

\begin{grammarlens}[title={What you've built}]
One more word for telling what happened.
\end{grammarlens}

\subsection*{Your turn}
\begin{itemize}
\raggedright
  \item \emph{Say it:} \emph{camatkār}
  \item \emph{Say it:} \emph{camatkār}, once more
  \item \emph{Say it:} \emph{badhāe kahyũ ke ā camatkār che.} — Everyone said it was a miracle.
  \item \emph{Recall:} say the Gujarati for a fairy, then the Gujarati for magic, then say \emph{camatkār} again
\end{itemize}

\begin{tcolorbox}[breakable,colback=teal!4,colframe=teal!35!black,title={Before you move on}]
What does \gu{ચમત્કાર} mean? (\textquotedblleft{}A miracle\textquotedblright{}.) Say it once more.
\end{tcolorbox}
\section[\texorpdfstring{\gu{ગાયબ} (gāyab)}{ગાયબ (gāyab)}]{\gu{ગાયબ} (gāyab) — vanished}

\label{lesson:GU-C333-gayab}



\begin{quote}
Before the new one: say the Gujarati for magic, then the Gujarati for a miracle.
\end{quote}

\subsection*{You'll want to know: \gu{ગાયબ}}
\textbf{\gu{ગાયબ}} — \emph{gāyab} — \textquotedblleft{}vanished\textquotedblright{}.

\begin{quote}
\gu{અચાનક} \gu{પરી} \gu{ગાયબ} \gu{થઈ} \gu{ગઈ}.
\end{quote}

\begin{quote}
\emph{acānak parī gāyab thaī gaī.}
\end{quote}

\begin{quote}
Suddenly the fairy vanished.
\end{quote}

\begin{cousinweb}
\textbf{\gu{ગાયબ}} (\emph{gāyab}) is usually traced, through Persian, to Arabic \emph{ġāʾib}, \textquotedblleft{}absent, hidden\textquotedblright{}. A thing or a person \textbf{\gu{ગાયબ} \gu{થાય} \gu{છે}}, disappears all at once.
\end{cousinweb}

\begin{grammarlens}[title={What you've built}]
One more word for telling what happened.
\end{grammarlens}

\subsection*{Your turn}
\begin{itemize}
\raggedright
  \item \emph{Say it:} \emph{gāyab}
  \item \emph{Say it:} \emph{gāyab}, once more
  \item \emph{Say it:} \emph{acānak parī gāyab thaī gaī.} — Suddenly the fairy vanished.
  \item \emph{Recall:} say the Gujarati for magic, then the Gujarati for a miracle, then say \emph{gāyab} again
\end{itemize}

\begin{tcolorbox}[breakable,colback=teal!4,colframe=teal!35!black,title={Before you move on}]
What does \gu{ગાયબ} mean? (\textquotedblleft{}Vanished\textquotedblright{}.) Say it once more.
\end{tcolorbox}
\section[\texorpdfstring{\gu{દેખાવું} (dekhāvũ)}{દેખાવું (dekhāvũ)}]{\gu{દેખાવું} (dekhāvũ) — to appear; to be seen}

\label{lesson:GU-C333-dekhavu}



\begin{quote}
Before the new one: say the Gujarati for a miracle, then the Gujarati for vanished.
\end{quote}

\subsection*{You'll want to know: \gu{દેખાવું}}
\textbf{\gu{દેખાવું}} — \emph{dekhāvũ} — \textquotedblleft{}to appear; to be seen\textquotedblright{}.

\begin{quote}
\gu{રાત્રે} \gu{બારીમાં} \gu{એક} \gu{પરી} \gu{દેખાઈ}.
\end{quote}

\begin{quote}
\emph{rātre bārīmā̃ ek parī dekhāī.}
\end{quote}

\begin{quote}
At night a fairy appeared at the window.
\end{quote}

Past (she …): \gu{તે} \gu{પરી} \gu{દેખાઈ}, \emph{te parī dekhāī}.

\begin{cousinweb}
\textbf{\gu{દેખાવું}} (\emph{dekhāvũ}) belongs with \emph{dekhvũ}, to see, which CDIAL lists under \emph{*dṛkṣati}, a form it reconstructs (CDIAL 6507): to be seen. What appears is the subject: \textbf{\gu{પરી} \gu{દેખાઈ}}, a fairy appeared.
\end{cousinweb}

\begin{grammarlens}[title={What you've built}]
One more word for telling what happened.
\end{grammarlens}

\subsection*{Your turn}
\begin{itemize}
\raggedright
  \item \emph{Say it:} \emph{dekhāvũ}
  \item \emph{Say it:} \emph{dekhāvũ}, once more
  \item \emph{Say it:} \emph{rātre bārīmā̃ ek parī dekhāī.} — At night a fairy appeared at the window.
  \item \emph{Recall:} say the Gujarati for a miracle, then the Gujarati for vanished, then say \emph{dekhāvũ} again
\end{itemize}

\begin{tcolorbox}[breakable,colback=teal!4,colframe=teal!35!black,title={Before you move on}]
What does \gu{દેખાવું} mean? (\textquotedblleft{}To appear; to be seen\textquotedblright{}.) Say it once more.
\end{tcolorbox}
